\documentclass{amsart}
% For dealing with references we use the comment environment
\usepackage{verbatim}
\newenvironment{reference}{\comment}{\endcomment}
%\newenvironment{reference}{}{}
\newenvironment{slogan}{\comment}{\endcomment}
\newenvironment{history}{\comment}{\endcomment}

% For commutative diagrams we use Xy-pic
\usepackage[all]{xy}

% We use 2cell for 2-commutative diagrams.
\xyoption{2cell}
\UseAllTwocells

% We use multicol for the list of chapters between chapters
\usepackage{multicol}

% This is generally recommended for better output
\usepackage{lmodern}
\usepackage[T1]{fontenc}

% For cross-file-references

% Package for hypertext links:
\usepackage{hyperref}

% For any local file, say "hello.tex" you want to link to please

% Theorem environments.
%
\theoremstyle{plain}
\newtheorem{theorem}[subsection]{Theorem}
\newtheorem{proposition}[subsection]{Proposition}
\newtheorem{lemma}[subsection]{Lemma}

\theoremstyle{definition}
\newtheorem{definition}[subsection]{Definition}
\newtheorem{example}[subsection]{Example}
\newtheorem{exercise}[subsection]{Exercise}
\newtheorem{situation}[subsection]{Situation}

\theoremstyle{remark}
\newtheorem{remark}[subsection]{Remark}
\newtheorem{remarks}[subsection]{Remarks}

\numberwithin{equation}{subsection}

% Macros
%
\def\lim{\mathop{\mathrm{lim}}\nolimits}
\def\colim{\mathop{\mathrm{colim}}\nolimits}
\def\Spec{\mathop{\mathrm{Spec}}}
\def\Hom{\mathop{\mathrm{Hom}}\nolimits}
\def\Ext{\mathop{\mathrm{Ext}}\nolimits}
\def\SheafHom{\mathop{\mathcal{H}\!\mathit{om}}\nolimits}
\def\SheafExt{\mathop{\mathcal{E}\!\mathit{xt}}\nolimits}
\def\Sch{\mathit{Sch}}
\def\Mor{\mathop{\mathrm{Mor}}\nolimits}
\def\Ob{\mathop{\mathrm{Ob}}\nolimits}
\def\Sh{\mathop{\mathit{Sh}}\nolimits}
\def\NL{\mathop{N\!L}\nolimits}
\def\CH{\mathop{\mathrm{CH}}\nolimits}
\def\proetale{{pro\text{-}\acute{e}tale}}
\def\etale{{\acute{e}tale}}
\def\QCoh{\mathit{QCoh}}
\def\Ker{\mathop{\mathrm{Ker}}}
\def\Im{\mathop{\mathrm{Im}}}
\def\Coker{\mathop{\mathrm{Coker}}}
\def\Coim{\mathop{\mathrm{Coim}}}

% Boxtimes
%
\DeclareMathSymbol{\boxtimes}{\mathbin}{AMSa}{"02}

%
% Macros for moduli stacks/spaces
%
\def\QCohstack{\mathcal{QC}\!\mathit{oh}}
\def\Cohstack{\mathcal{C}\!\mathit{oh}}
\def\Spacesstack{\mathcal{S}\!\mathit{paces}}
\def\Quotfunctor{\mathrm{Quot}}
\def\Hilbfunctor{\mathrm{Hilb}}
\def\Curvesstack{\mathcal{C}\!\mathit{urves}}
\def\Polarizedstack{\mathcal{P}\!\mathit{olarized}}
\def\Complexesstack{\mathcal{C}\!\mathit{omplexes}}
% \Pic is the operator that assigns to X its picard group, usage \Pic(X)
% \Picardstack_{X/B} denotes the Picard stack of X over B
% \Picardfunctor_{X/B} denotes the Picard functor of X over B
\def\Pic{\mathop{\mathrm{Pic}}\nolimits}
\def\Picardstack{\mathcal{P}\!\mathit{ic}}
\def\Picardfunctor{\mathrm{Pic}}
\def\Deformationcategory{\mathcal{D}\!\mathit{ef}}

\setlength{\emergencystretch}{3em}
\begin{document}
\title[Stacks Project, Tag 0BDC]{577 --- Existence of an ample line bundle descends from an arbitrary field extension}
\maketitle
\noindent Copyright (C) 2005--2025 Johan de Jong. Stacks Project authors. Source: \href{https://stacks.math.columbia.edu/tag/0BDC}{Tag 0BDC}.

\noindent Editorial difficulty note (not author text). Difficulty H1: The proof first recovers quasi-compactness and quasi-separatedness from the extended scheme. It descends the line bundle and its ampleness to a finite-type ground-field algebra, specializes to a finite residue extension and applies finite locally free ample descent..

\noindent Editorial provenance note (not author text). History: excerpt from revision \texttt{a0444\allowbreak 6e57e\allowbreak c1fbc\allowbreak 252a8\allowbreak 71afc\allowbreak ec775\allowbreak 2fb28\allowbreak 07b14}, varieties.tex, lines 2739--2780. Reference presentation was adapted; original-author context and core support blocks were retained or added. The mathematical wording is unchanged. Licensed under GNU FDL 1.2 or later; no invariant sections or cover texts. See ../licenses/COPYING. Context blocks reproduce author definitions and supporting results; they are not additional counted problems.

\begin{definition}
\label{properties-definition-ample}
\begin{reference}
\href{https://stacks.math.columbia.edu/bibliography}{Bibliography: EGA (II Definition 4.5.3)}
\end{reference}
Let $X$ be a scheme.
Let $\mathcal{L}$ be an invertible $\mathcal{O}_X$-module.
We say $\mathcal{L}$ is {\it ample} if
\begin{enumerate}
\item $X$ is quasi-compact, and
\item for every $x \in X$ there exists an $n \geq 1$
and $s \in \Gamma(X, \mathcal{L}^{\otimes n})$ such
that $x \in X_s$ and $X_s$ is affine.
\end{enumerate}
\end{definition}

\begin{lemma}
\label{lemma-ample-after-field-extension}
Let $k$ be a field. Let $X$ be a scheme over $k$.
If there exists an ample invertible sheaf on $X_K$ for some
field extension $K/k$, then $X$ has an ample invertible
sheaf.
\end{lemma}

\begin{proof}
Let $K/k$ be a field extension such that $X_K$ has
an ample invertible sheaf $\mathcal{L}$.
The morphism $X_K \to X$ is surjective. Hence $X$ is quasi-compact
as the image of a quasi-compact scheme (Properties, Definition
\ref{properties-definition-ample}). Since $X_K$ is quasi-separated
(by Properties, Lemma
\href{https://stacks.math.columbia.edu/tag/01PY}{lemma affine s opens cover quasi separated (Tag 01PY)})
we see that $X$ is quasi-separated: If $U, V \subset X$ are
affine open, then $(U \cap V)_K = U_K \cap V_K$ is quasi-compact
and $(U \cap V)_K \to U \cap V$ is surjective. Thus
Schemes, Lemma \href{https://stacks.math.columbia.edu/tag/01KO}{lemma characterize quasi separated (Tag 01KO)} applies.

\medskip\noindent
Write $K = \colim A_i$ as the colimit of the subalgebras of $K$
which are of finite type over $k$. Denote
$X_i = X \times_{\Spec(k)} \Spec(A_i)$.
Since $X_K = \lim X_i$ we find an $i$ and an invertible sheaf
$\mathcal{L}_i$ on $X_i$ whose pullback to $X_K$ is $\mathcal{L}$
(Limits, Lemma \href{https://stacks.math.columbia.edu/tag/0B8W}{lemma descend invertible modules (Tag 0B8W)};
here and below we use that $X$ is quasi-compact and quasi-separated as
just shown). By Limits, Lemma \href{https://stacks.math.columbia.edu/tag/09MT}{lemma limit ample (Tag 09MT)}
we may assume $\mathcal{L}_i$ is ample after possibly increasing $i$.
Fix such an $i$ and let $\mathfrak m \subset A_i$ be a maximal
ideal. By the Hilbert Nullstellensatz
(Algebra, Theorem \href{https://stacks.math.columbia.edu/tag/00FV}{theorem nullstellensatz (Tag 00FV)})
the residue field $k' = A_i/\mathfrak m$ is a finite
extension of $k$. Hence $X_{k'} \subset X_i$ is a closed subscheme
hence has an ample invertible sheaf
(Properties, Lemma \href{https://stacks.math.columbia.edu/tag/01PU}{lemma ample on closed (Tag 01PU)}).
Since $X_{k'} \to X$ is finite locally free we conclude
that $X$ has an ample invertible sheaf by
Divisors, Proposition \href{https://stacks.math.columbia.edu/tag/0BD4}{proposition push down ample (Tag 0BD4)}.
\end{proof}
\end{document}
